% =====================================================================
\chapter{SOLUTION DEVELOPMENT}
\label{ch:solution}
% =====================================================================

This chapter develops the solution selected in Chapter~\ref{ch:method}.
Section~\ref{sec:model} identifies the components, parameters, assumptions
and constraints of the system and formulates the mathematical model of the
mechanism. Section~\ref{sec:development} develops the solution: the overall
procedure, the route-generation algorithm, the local pruning frontier and its
worst-case size, the label rule, payment computation, a complete small
numerical solution, the implementation, the investigation of alternative
solutions, and a comparison of alternatives under feasibility constraints and
standards.

\section{Model / Prototype Solution}
\label{sec:model}

\subsection{Components and parameters of the system}

A dispatch session is static and deterministic. It consists of a finite set
of orders $O$ and a set of strategic drivers $N=G\cup C$, where $G$ is the
set of gigworkers and $C$ the set of occasional drivers.

\paragraph{Orders.} Order $o\in O$ has a pickup node $p_o$, a delivery node
$d_o$, a demand, a release time, and time windows at both nodes; an arrival
before the opening of a window is met by waiting. Every order can be
delegated to the fixed-capacity fleet (FD) at a public price $q_o\ge0$; for
$A\subseteq O$ we write $q(A)=\sum_{o\in A}q_o$.

\paragraph{Network.} Travel distances $d(\cdot,\cdot)$ and times
$\tau(\cdot,\cdot)$ are nonnegative and satisfy the triangle inequality, and
service times are nonnegative.

\paragraph{Drivers.} A gigworker has a start location, an availability
interval and a capacity; her route is open and ends at its last delivery. An
occasional driver has an origin, a personal destination, a direct trip from
one to the other, a detour budget $\tau_i$, an availability interval and a
capacity; her route ends at the destination and its detour time may not
exceed $\tau_i$. These data are public or committed before the session
opens. Each driver also has a public distance coefficient $\kappa_i$ and a
private value of time $\theta_i$.

\paragraph{Platform parameters.} The platform publishes a bundle cap $B$ and,
when pruning is used, a report domain $\Theta=[\lth,\uth]$.
Table~\ref{tab:params} lists the parameter values used in the computational
study.

\begin{table}[t]
\centering
\caption{Parameter values of the computational study}
\label{tab:params}
\small
\begin{tabular}{P{6.0cm}P{8.4cm}}
\toprule
Parameter & Value\\
\midrule
Number of orders $n$ & 10, 15, 20 (main grid); 25, 30 (robustness); up to 100 (route-generation scaling study); 50, 75, 100 and beyond (planned large-scale experiments)\\
Supply mix $(|G|,|C|)$ & (2,2), (3,2), (2,3), (3,3); one more of each at $n\ge 25$\\
Bundle cap $B$ & 3; 4 as a sensitivity case for $n\le15$\\
Time-window width & 120 min (60 and 240 in robustness tests)\\
Detour budget $\tau_i$ & 30 min (20 and 45 in robustness tests)\\
Service time per stop; speed & 5 min; 20 km/h\\
Distance coefficient $\kappa_i$ & 1 cost unit per km\\
Value of time $\theta_i$ & $U[18,25]$ cost units per hour ($U[15,30]$ in one robustness test)\\
FD price $q_o$ & $8+3\max\{0,\mathrm{dist}_o-1\}$, $\mathrm{dist}_o$ in km (scaled by 0.75 and 1.25 in robustness tests)\\
Report domain $\Theta$ & $[18,25]$\\
Alignment of OD trips with demand & 0.10, 0.30, 0.50, 0.70, 0.90\\
Replications & 25 per cell (1,500 instances in the main grid)\\
\bottomrule
\end{tabular}
\end{table}

\subsection{Assumptions and constraints}

A route $r$ of driver $i$ serves a bundle $S_r\subseteq O$ with
$1\le|S_r|\le B$. It is \emph{feasible} if:
\begin{enumerate}[label=(F\arabic*),leftmargin=1.3cm,itemsep=1pt]
\item the pickup of every order precedes its delivery, and both are served
by the same driver;
\item every order has been released when it is picked up;
\item service at every stop starts within its time window, waiting if the
driver arrives early;
\item the load after every stop lies between zero and the driver's capacity;
\item the route is completed within the driver's availability interval;
\item for an occasional driver, the route ends at her destination and its
detour time does not exceed $\tau_i$.
\end{enumerate}
The assumptions of Section~\ref{sec:scope} apply. In addition, the
mechanism itself imposes three constraints: every order is served exactly
once, by a route or by FD; every driver receives at most one route; and the
pools $R_i$ are fixed before any report is collected.

\subsection{Cost model}
\label{sec:sd-costs}

Serving route $r$ costs driver $i$
\begin{equation}
  C_{ir}(\theta_i)=K_{ir}+\theta_i\,W_{ir},
  \label{eq:cost}
\end{equation}
where $K_{ir}\ge 0$ is a distance cost, $W_{ir}\ge 0$ is working time in
hours, and $\theta_i>0$ is the driver's private value of time. Both
coefficients are computed by the platform from public data. For a gigworker
with departure time $t_0$,
\[
K_{ir}=\kappa_i\cdot(\text{route length}),\qquad
W_{ir}=(\text{completion time}-t_0)/60,
\]
with times in minutes. For an occasional driver, costs are incremental over
the direct trip:
\begin{align*}
K_{ir}&=\kappa_i\max\{0,\text{route length}-\text{direct length}\},\\
W_{ir}&=\max\{0,\text{route time}-\text{direct time}\}/60 .
\end{align*}
A driver who reports $b_i$ in place of $\theta_i$ is evaluated at the
\emph{reported cost} $c_{ir}(b_i)=K_{ir}+b_iW_{ir}$. The idle option of
driver $i$ is the empty route $0_i$ with $S_{0_i}=\emptyset$ and $K=W=0$.

\paragraph{Why one private parameter suffices.} A driver's true cost has two
dimensions, a cost per kilometre and a cost per hour, and the single-parameter
form~\eqref{eq:cost} treats the first as public. This is an empirical
approximation, not merely a convenience. In last-mile geometry route time is
close to route length divided by speed plus service time, so the vectors
(length, time) of a driver's bundle menu lie close to one ray. In
preliminary measurements on instances of the same generator, the angular
spread of a driver's menu was about one degree, the median number of distinct
bundle orderings over the type range was one under tight time windows, and
the welfare loss from projecting a two-dimensional cost onto one parameter
was 0.002--0.05\%. The conclusion was unchanged when the variability of
speed was inflated eightfold: a change of speed rotates a driver's menu
rather than spreading it, and the spread is governed by the ratio of service
time to route length.

\subsection{Winner-determination model}

Let $R_i$ be the route pool of driver $i$, fixed before reports are
collected. Given a report profile $b$, the winner-determination problem
(Algorithm~B) is the set-partitioning model
\begin{align}
V(R;b)=\min_{x,z}\;&\sum_{i\in N}\sum_{r\in R_i}c_{ir}(b_i)\,x_{ir}+\sum_{o\in O}q_o z_o
\label{eq:wdp}\\
\text{s.t.}\;&\sum_{i\in N}\sum_{r\in R_i:\,o\in S_r}x_{ir}+z_o=1 && o\in O,\label{eq:cover}\\
&\sum_{r\in R_i}x_{ir}\le 1 && i\in N,\label{eq:unit}\\
&x_{ir},z_o\in\{0,1\} && i\in N,\ r\in R_i,\ o\in O,\notag
\end{align}
where $x_{ir}=1$ if driver $i$ receives route $r$ and $z_o=1$ if order $o$
is delegated to FD. The objective~\eqref{eq:wdp} is the total reported cost
of the selected routes plus the FD cost of the delegated orders.
Constraints~\eqref{eq:cover} ensure that every order is served exactly once,
and constraints~\eqref{eq:unit} that every driver receives at most one route.
The problem is always feasible, because FD is uncapacitated: setting every
$z_o=1$ is a solution of value $q(O)$.

Ties are broken deterministically and independently of reports: among
optimal solutions the model minimises the number of orders sent to FD and
then chooses lexicographically by route identifier. $V_{-k}(R;b)$ denotes the
optimal value of~\eqref{eq:wdp} with driver $k$ removed.

\subsection{Payment rule and truthfulness}

Algorithm~C pays each winner $i$ with selected route $r_i^\ast$
\begin{equation}
p_i=c_{ir_i^\ast}(b_i)+V_{-i}(R;b)-V(R;b),
\label{eq:payment}
\end{equation}
and pays losers nothing. Only solutions proven optimal with zero gap are
used.

\begin{proposition}[Truthfulness; adapted from the Groves--VCG argument]
\label{prop:dsic}
If the pools $R_i$ are fixed before reports are observed, every base and
removal problem is solved exactly with report-independent tie-breaking, and
each report is a nonnegative scalar entering the objective as in
\eqref{eq:wdp}, then truthful reporting is a dominant strategy for every
driver and every winner's utility is nonnegative.
\end{proposition}

This is the standard maximal-in-range argument
\citep{nisan2007computationally}: driver $i$'s utility under
report $b_i$ equals $V_{-i}$, which she cannot influence, minus the
\emph{true} total cost of the allocation chosen at $b_i$; truthful reporting
makes the mechanism minimise exactly that true total cost over the fixed
range. The argument does not depend on whether a route belongs to a
gigworker or an occasional driver, and it is not claimed as a contribution.

\subsection{Report domain}
\label{sec:reportdomain}

The basic mechanism places no restriction on reports. The pruning results of
Sections~\ref{sec:sd-frontier} and~\ref{sec:sd-label} are stated for a report
domain $\Theta=[\lth,\uth]$ with $0\le\lth<\uth$, announced before bidding;
pruning is exact for every report profile in $\Theta^N$. When pruning is
used, the message space is therefore restricted to $\Theta$. Setting
$\Theta=[0,\infty)$ removes the restriction at the price of a larger frontier
(Remark~\ref{rem:monotone}). In the experiments $\Theta=[18,25]$, which
coincides with the support of the baseline type distribution.

\section{Solution Development}
\label{sec:development}

\subsection{Overview of the procedure}

Procedure~\ref{alg:mechanism} states the complete mechanism. Lines 1--4 use
public data only; the frontier step is optional and, by
Theorem~\ref{thm:frontier}, does not change any outcome.

\begin{algorithm}[!htbp]
\caption{Truthful route-bundle procurement}
\label{alg:mechanism}
\begin{algorithmic}[1]
\Require orders $O$, FD prices $q$, bundle cap $B$, report domain $\Theta$, public driver data
\Ensure allocation $(x^\ast,z^\ast)$ and payments $(p_i)_{i\in N}$
\For{each driver $i\in N$} \Comment{no report is read in lines 1--4}
  \State $R_i\gets\textsc{GenerateRoutes}(i)$ \Comment{Algorithm A, Proc.~\ref{alg:routegen}}
  \State $R_i\gets\textsc{Frontier}(R_i,q,\Theta)$ \Comment{optional, Proc.~\ref{alg:frontier}}
\EndFor
\State publish a hash of $(R_i)_{i\in N}$; collect one report $b_i\in\Theta$ per driver
\State solve~\eqref{eq:wdp} on $R$ at $b$ to optimality $\to(x^\ast,z^\ast)$, $V(R;b)$ \Comment{Algorithm B}
\State $(p_i)_{i\in N}\gets\textsc{Payments}(R,b,x^\ast)$ \Comment{Algorithm C, Proc.~\ref{alg:payments}}
\State \Return $(x^\ast,z^\ast)$, $(p_i)_{i\in N}$
\end{algorithmic}
\end{algorithm}

% ---------------------------------------------------------------------
\subsection{Algorithm A: bid-independent route generation}
\label{sec:sd-algA}

Algorithm~A enumerates, for each driver separately, all feasible routes
with at most $B$ orders and keeps, for each bundle, the Pareto-minimal
pairs $(K,W)$. It never reads a report.

\paragraph{Labels.} A label is a tuple $\ell=(v,IV,C,t,K,W)$ consisting of
the current node $v$, the set $IV$ of orders on board, the set $C$ of
delivered orders, the clock $t$ after service at $v$, and the accumulated
distance and time costs. A pickup of order $j$ is admissible if
$j\notin IV\cup C$, $|IV|+|C|<B$ and the pickup is reached before the
closing of its window; a delivery requires $j\in IV$; an occasional driver
may return to her destination once $IV=\emptyset$, provided the detour budget
is respected. Each transition waits for the opening of the next window if
necessary. A gigworker's label is complete when $IV=\emptyset$ and
$C\neq\emptyset$; an occasional driver's label is complete when it has
reached the destination, after which the direct trip is subtracted once.
Labels are processed in rounds of increasing $|IV|+|C|$.

\paragraph{Dominance.} Two labels are compared only if they share the key
$(v,IV,C)$; $\ell_A$ dominates $\ell_B$ if $t_A\le t_B$, $K_A\le K_B$ and
$W_A\le W_B$ with at least one strict inequality. This is the forward
dominance of \citet{dumas1991pickup}. Labels with the same key have exactly
the same feasible continuations, so every completion of $\ell_B$ is matched
by a no-worse completion of $\ell_A$. The key cannot be coarsened in
general. If two labels at the same node with the same orders on board have
delivered different orders, and an order delivered by one of them can no
longer be picked up in time by the other, then their completions serve
different bundles. Merging labels under the coarser key $(v,IV)$ can
therefore discard the only label able to reach some bundle. On 720 instances with $n\le 6$, both
coarsenings that were tested (dropping $C$, or keeping only $|C|$) disagreed
with brute-force enumeration in 46--47\% of instances.

\paragraph{Pruning layers.} Procedure~\ref{alg:routegen} states Algorithm~A
in full. It contains three pruning layers. Layer~1 is the label dominance
above. Layer~2 keeps, for each bundle, the Pareto-minimal pairs $(K,W)$ of
the completed routes; it never compares routes with different bundles,
because the orders missing from the smaller bundle would have to be served
by someone. Layer~3 is the optional FD-completion label rule of
Section~\ref{sec:sd-label}. Table~\ref{tab:layers} contrasts the three
layers with the frontier filter of Section~\ref{sec:sd-frontier}.

\begin{algorithm}[!htbp]
\caption{\textsc{GenerateRoutes}$(i)$: Algorithm A for one driver}
\label{alg:routegen}
\begin{algorithmic}[1]
\Require public data of driver $i$; orders $O$; cap $B$; flag \textit{useRule}
\Ensure route pool $R_i$ (Pareto-minimal $(K,W)$ per bundle)
\State $\mathcal L_0\gets\{(v_0,\emptyset,\emptyset,t_0,0,0)\}$; $\mathcal L_k\gets\emptyset$ for $k=1,\dots,B$; $\mathcal C\gets\emptyset$
\For{$k\gets 0$ \textbf{to} $B$} \Comment{$k=|IV|+|C|$, the number of touched orders}
  \State $\mathcal Q\gets\mathcal L_k$
  \While{$\mathcal Q\ne\emptyset$}
    \State $\mathcal Q\gets\textsc{FilterDominated}(\mathcal Q)$ \Comment{Layer 1: same key $(v,IV,C)$}
    \State $\mathcal Q'\gets\emptyset$
    \For{each label $\ell\in\mathcal Q$}
      \If{$i\in G$ \textbf{and} $IV(\ell)=\emptyset$ \textbf{and} $C(\ell)\ne\emptyset$}
        \State add $\ell$ to $\mathcal C$ \Comment{complete gigworker route; $\ell$ may still be extended}
      \EndIf
      \If{$i\in C$ \textbf{and} $IV(\ell)=\emptyset$ \textbf{and} $C(\ell)\ne\emptyset$}
        \State $\ell'\gets\textsc{TryHome}(\ell)$; \textbf{if} $\ell'\ne\bot$ \textbf{then} add $\ell'$ to $\mathcal C$ \Comment{checks the detour budget}
      \EndIf
      \For{each $j\in IV(\ell)$}
        \State $\ell'\gets\textsc{TryDelivery}(\ell,j)$
        \If{$\ell'\ne\bot$ \textbf{and not} (\textit{useRule} \textbf{and} $\textsc{RuleFires}(\ell')$)} add $\ell'$ to $\mathcal Q'$ \EndIf
      \EndFor
      \If{$k<B$}
        \For{each $j\in O\setminus(IV(\ell)\cup C(\ell))$}
          \State $\ell'\gets\textsc{TryPickup}(\ell,j)$
          \If{$\ell'\ne\bot$ \textbf{and not} (\textit{useRule} \textbf{and} $\textsc{RuleFires}(\ell')$)} add $\ell'$ to $\mathcal L_{k+1}$ \EndIf
        \EndFor
      \EndIf
    \EndFor
    \State $\mathcal Q\gets\mathcal Q'$ \Comment{each delivery adds one order to $C$}
  \EndWhile
\EndFor
\State for occasional drivers, subtract the direct trip from each complete label
\For{each bundle $S$ appearing in $\mathcal C$}
  \State $R_i(S)\gets\textsc{ParetoFront}\bigl(\{(K_\ell,W_\ell):\ell\in\mathcal C,\ C(\ell)=S\}\bigr)$ \Comment{Layer 2}
\EndFor
\State \Return $R_i=\bigcup_S R_i(S)$, one route per distinct pair $(K,W)$
\end{algorithmic}
\end{algorithm}

\begin{table}[t]
\centering
\caption{Pruning layers of Algorithm A and the frontier filter}
\label{tab:layers}
\small
\begin{tabular}{P{3.1cm}P{3.0cm}P{3.2cm}P{1.5cm}P{2.9cm}}
\toprule
Filter & Compares & Scope & Uses $q$? & When\\
\midrule
Layer 1: label dominance & label vs.\ label & same $(v,IV,C)$ & No & During generation\\
Layer 2: per-bundle Pareto & route vs.\ route & same bundle $S$ & No & After generation\\
Layer 3: label rule (Sec.~\ref{sec:sd-label}) & partial route vs.\ FD, lower bound & any $T\subseteq P(L)$ & Yes & During generation\\
Frontier $\Pi^\star$ (Sec.~\ref{sec:sd-frontier}) & route vs.\ subset route $+$ FD, exact & any $r'$ with $S_{r'}\subseteq S_r$ & Yes & After generation\\
\bottomrule
\end{tabular}
\end{table}

\paragraph{Properties.} Algorithm~A is complete: for $n\le 6$ its output
coincides with brute-force enumeration. For fixed $B$ the pool of a driver
has size $O(|O|^B(2B)!/2^B)$. The pool, and a hash of it, are invariant
under any change of reports. Keeping only same-bundle Pareto-minimal pairs
is safe because $b\ge 0$: if $K_1\le K_2$ and $W_1\le W_2$, then
$K_1+bW_1\le K_2+bW_2$ for every report.

\paragraph{Scaling.} Before the main experiments, Algorithm~A was run on
instances with up to $n=100$ orders in a $20\times20$\,km area, for bundle caps
$B\in\{2,3,4\}$ and time-window widths from 30 to 240 minutes. Its
correctness on the same generator was first checked against brute-force
enumeration on 180 small cases ($n\le 7$, no violation).
Table~\ref{tab:scaling} reports running times. Grids with $n\le 50$ were
run in full (mean over up to 20 instance--driver-count combinations); the
cells for $n=75$ and $n=100$ come from a probe with one driver per point and
are therefore per driver. With $B=2$ the algorithm stays below one second up
to $n=100$. With $B=3$ it handles $n=75$ for every window width, and $n=100$
when windows are at most 120 minutes. With $B=4$ it exceeds the 90-second
probe limit before $n=50$ for most window widths. The bundle cap, not the
number of orders, is therefore the main driver of cost, consistent with the
$O(|O|^B)$ bound, and large instances must use $B\le 3$.

\begin{table}[t]
\centering
\caption{Running time of Algorithm A against the number of orders (selected cells)}
\label{tab:scaling}
\small
\begin{tabular}{cccccccc}
\toprule
$B$ & Window (min) & $n=10$ & $n=20$ & $n=30$ & $n=50$ & $n=75$\textsuperscript{a} & $n=100$\textsuperscript{a}\\
\midrule
2 & 120 & 0.01 s & 0.02 s & 0.06 s & 0.31 s & 0.3 s & 0.5 s\\
2 & 240 & 0.01 s & 0.03 s & 0.09 s & 0.44 s & 0.5 s & 0.7 s\\
3 & 60  & 0.03 s & 0.18 s & 0.78 s & 8.07 s & 14 s & 33 s\\
3 & 120 & 0.05 s & 0.35 s & 1.69 s & 16.9 s & 30 s & 77 s\\
3 & 240 & 0.15 s & 1.02 s & 4.99 s & 46.2 s & 76 s & $>90$ s\\
4 & 60  & 0.13 s & 2.29 s & 14.9 s & 218 s & -- & --\\
4 & 120 & 0.33 s & 7.81 s & 44.0 s & timeout & -- & --\\
\bottomrule
\end{tabular}
\par\smallskip\footnotesize \textsuperscript{a}Single-driver probe, time per driver; other columns are means over the instances of a grid cell. ``$>90$ s'' exceeded the probe limit; ``--'' not run.
\end{table}

\paragraph{The asymmetry between the classes.} For an occasional driver the
detour budget provides a monotone bound: a lower bound on the detour of a
bundle eliminates the bundle and all its supersets before any sequence is
enumerated \citep{li2026auction}. For a gigworker the only constraints are
time windows. Two orders may each be servable while the pair is not, and a
bundle of distant orders remains feasible whenever the windows are wide, so
no bound as strong as a detour budget is available. On the experimental
instances the pool of an occasional driver is 0.7--1.7\% of that of a
gigworker. Reducing the gigworkers' pools safely is therefore the purpose of
the frontier of Section~\ref{sec:sd-frontier}.

% ---------------------------------------------------------------------
\subsection{The local pruning frontier}
\label{sec:sd-frontier}

When the platform considers deleting a route $r$ of driver $i$ without
looking at the other drivers, the only substitutes it can certify are those
that involve nobody else: a route of the same driver that serves part of
$S_r$, with the remainder sent to FD. The frontier keeps exactly the routes
for which no such substitute is at least as cheap for every admissible
report. This section proves that deleting all other routes is safe: it
changes neither the optimal value, nor any counterfactual value, nor any
payment.

\begin{definition}[Local view, local rule, safety]
\label{def:local}
The local view of driver $i$ is $L_i=(O,q,B,\Theta,R_i)$. A local pruning
rule $\Pi$ is a deterministic map that assigns to every local view a set
$\Pi(L_i)\subseteq R_i$ of routes to delete; it is applied to every driver
and yields pools $R^\Pi_j=R_j\setminus\Pi(L_j)$. An instance is any finite
set of drivers whose pools consist of routes with bundles of size at most
$B$ and nonnegative coefficients. The rule is safe if $V(R^\Pi;b)=V(R;b)$
for every instance and every $b\in\Theta^N$.
\end{definition}

A local rule reads no report, so it leaves the allocation range fixed before
bidding and Proposition~\ref{prop:dsic} continues to apply to the pruned
mechanism.

\begin{definition}[FD-completion envelope and local frontier]
\label{def:frontier}
For $r\in R_i$ let $D(r)=\{r'\in R_i\cup\{0_i\}\setminus\{r\}:S_{r'}\subseteq S_r\}$
and
\begin{equation}
E_r(b)=\min_{r'\in D(r)}\bigl\{c_{r'}(b)+q(S_r\setminus S_{r'})\bigr\},\qquad
\mu_r(b)=E_r(b)-c_r(b).
\label{eq:envelope}
\end{equation}
The local frontier of driver $i$ is
$\Kstar_i=\{r\in R_i:\max_{b\in\Theta}\mu_r(b)>0\}$, and the rule $\Pi^\star$
deletes $R_i\setminus\Kstar_i$ for every driver.
\end{definition}

In words, $E_r(b)$ is the cheapest way to serve the bundle of $r$ using a
\emph{smaller} route of the same driver and sending the missing orders to
FD, and $\mu_r(b)$ is the margin by which $r$ beats that alternative. A route
is kept if its margin is positive for at least one admissible report.
Because $0_i\in D(r)$, $E_r(b)\le q(S_r)$. The envelope $E_r$ is the minimum
of finitely many affine functions of $b$, hence concave and piecewise linear,
so $\mu_r$ is concave and its maximum over $\Theta$ is attained at an
endpoint or at a breakpoint of $E_r$. One assumption on the pools is used.

\begin{assumption}[No duplicates]\label{ass:nodup}
Distinct routes of $R_i$ with the same bundle have distinct pairs $(K,W)$.
\end{assumption}

Algorithm~A satisfies Assumption~\ref{ass:nodup} by construction, because
Layer~2 returns one route per distinct pair $(K,W)$ of each bundle.

\paragraph{An example.} The following example shows how the frontier
decides, and why a cruder rule that compares routes with different bundles
is not safe.

\begin{example}[Twin instances]
\label{ex:twin}
Let gigworker $g$ face orders $o_1,o_2$ with $q_{o_1}=15$, $q_{o_2}=45$ and
$\Theta=[18,25]$, and own three routes: $r_a$ serving $\{o_1\}$ with
$(K,W)=(6,0.6)$; $r_b$ serving $\{o_2\}$ with $(8,0.8)$; and $r_c$ serving
$\{o_1,o_2\}$ with $(11,1.2)$. Route $r_a$ is outside the frontier, since
$E_{r_a}=q_{o_1}=15<6+0.6b$ on $\Theta$. Route $r_b$ is in the frontier,
since $E_{r_b}=45>8+0.8b$. Route $r_c$ is in the frontier, since
$E_{r_c}(18)=\min\{6+10.8+45,\;8+14.4+15,\;60\}=37.4>32.6=c_{r_c}(18)$. In
instance $I_0$, in which $g$ is the only driver, $r_c$ is optimal on all of
$\Theta$ and $r_b$ is never optimal. In instance $I_1$, which adds a driver
who serves $o_1$ at cost $0.5$, the optimum at $b=18$ is $22.4+0.5=22.9$ and
uses $r_b$; without $r_b$ it is $32.6$. A rule that deletes $r_b$ --- for
instance a per-driver hull across bundles, under which $r_a$ dominates $r_b$
componentwise --- is correct on $I_0$ and wrong on $I_1$, and no rule that
sees only $g$'s data can tell the two instances apart.
\end{example}

Figure~\ref{fig:frontier-intuition} draws the three margins of
Example~\ref{ex:twin}.

\begin{figure}[t]
\centering
\includegraphics[width=\textwidth]{fig_c1_frontier_intuition.pdf}
\caption{Local margins in Example~\ref{ex:twin}. Solid lines: reported cost
$c_r(b)$; dashed: FD-completion envelope $E_r(b)$ (dotted in (c): its three
affine branches); grey band: report domain $\Theta=[18,25]$. A route belongs
to $\Kstar$ if and only if the envelope lies strictly above its cost
somewhere in $\Theta$ (shaded). Route $r_a$ is deleted only because $\Theta$
excludes reports below 15; with unrestricted reports it would be retained
(Remark~\ref{rem:monotone}).}
\label{fig:frontier-intuition}
\end{figure}

\paragraph{Main result.} The next theorem answers TQ1: deleting every route
outside the frontier is safe.

\begin{theorem}[Safety of the local frontier]
\label{thm:frontier}
Let $\lth<\uth$ and let every pool satisfy Assumption~\ref{ass:nodup}. Then
$\Pi^\star$ is safe: $V(R^{\Pi^\star};b)=V(R;b)$ for every instance and every
$b\in\Theta^N$. Moreover $V_{-k}(R^{\Pi^\star};b)=V_{-k}(R;b)$ for every
driver $k$ and every $b\in\Theta^N$.
\end{theorem}

The difficulty is that a deleted route must be replaced by a route that is
itself \emph{retained}. Its envelope may be attained by a subset route that
is also outside the frontier, which would then also be deleted. The
following lemma shows that a retained substitute always exists.

\begin{lemma}[A retained substitute exists]
\label{lem:substitute}
Let $R_i$ satisfy Assumption~\ref{ass:nodup} and $\lth<\uth$. For every
$r\in R_i\setminus\Kstar_i$ and every $b\in\Theta$ there is
$u\in\Kstar_i\cup\{0_i\}$ with $S_u\subseteq S_r$ and
$c_u(b)+q(S_r\setminus S_u)\le c_r(b)$.
\end{lemma}

\paragraph{Idea of the proof.} Fix $b$ and consider all ways of serving
$S_r$ with $r$ itself or with one subset route plus FD. Among the cheapest of
these, take one whose bundle $T$ is as small as possible. If $T$ is empty,
the idle route $0_i$ is the substitute. Otherwise every smaller alternative is
strictly more expensive at $b$, and these strict inequalities between lines
in the report remain true for reports close to $b$. Among the cheapest routes
with bundle $T$, which have pairwise different slopes $W$ by
Assumption~\ref{ass:nodup}, the one with the smallest slope (or the largest,
if $b=\uth$) becomes strictly cheapest after a small move of the report
inside $\Theta$. At that report its margin is positive, so it belongs to
$\Kstar_i$. The full proof is in Appendix~\ref{app:proofs}.

\begin{proof}[Proof of Theorem~\ref{thm:frontier}]
Fix an instance, a profile $b\in\Theta^N$ and an optimal allocation $\rho$
for the full pools. For every driver $j$ whose route $\rho_j$ is not in
$\Kstar_j\cup\{0_j\}$, Lemma~\ref{lem:substitute} gives
$u_j\in\Kstar_j\cup\{0_j\}$ with $S_{u_j}\subseteq S_{\rho_j}$ and
$c_{u_j}(b_j)+q(S_{\rho_j}\setminus S_{u_j})\le c_{\rho_j}(b_j)$. Replace
$\rho_j$ by $u_j$ and send $S_{\rho_j}\setminus S_{u_j}$ to FD. Bundles only
shrink, so they remain disjoint; every order is still served once; and no
other driver is affected, so all replacements can be made simultaneously.
The resulting allocation uses only pruned pools and costs at most
$V(R;b)$, so $V(R^{\Pi^\star};b)\le V(R;b)$; the reverse inequality holds
because pruned pools are subsets. Removing driver $k$ yields another
instance in which every remaining driver has the same local view, so the
same argument applies to $V_{-k}$.
\end{proof}

\begin{corollary}[VCG outcomes are preserved]
\label{cor:vcg}
Running Algorithms~B and~C on $R^{\Pi^\star}$ is VCG on a bid-independent
range and is therefore truthful and individually rational. At every profile
at which the full problem has a unique optimal allocation, the pruned problem
has the same allocation and every payment~\eqref{eq:payment} coincides with
its full-pool value.
\end{corollary}

\begin{proof}
The first statement is Proposition~\ref{prop:dsic}. Under uniqueness, any
pruned optimum has value $V(R;b)$ by Theorem~\ref{thm:frontier} and is
feasible for the full pools, hence equal to the full optimum; $V$ and every
$V_{-k}$ coincide by Theorem~\ref{thm:frontier}.
\end{proof}

\paragraph{Scope of the result.}

\begin{remark}[What the frontier does not do]
\label{rem:scope}
The theorem guarantees safety, not that the frontier is small. Rules outside
the local class can keep fewer routes. (i) \emph{Non-local rules} that use
the public data of other drivers remain bid-independent, but deciding
exactly which routes they may delete requires reasoning about the other
drivers' winner-determination problem, which is NP-hard for $B\ge 3$ already
at a single report profile: with one driver per 3-set, each owning one route
of cost zero, and $q_o=1$, the optimal value is zero if and only if an exact
cover by 3-sets exists. (ii) \emph{Instance-specific rules}, which only need
to be correct on the instance at hand, can retain far less: on five
experimental instances only 0.4--1.5\% of generated routes are optimal for
any of 1,000 sampled report profiles, whereas the frontier retains 3--23\% of
the pool on the same instances.
\end{remark}

\begin{remark}[Monotonicity]
\label{rem:monotone}
$E_r$ is nondecreasing in every $q_o$, and the maximum of $\mu_r$ over a
larger interval is larger. The frontier therefore grows with the FD tariff
and with the width of $\Theta$: local pruning is strong when FD is cheap and
weak when it is expensive. With $\Theta=[0,\infty)$ the rule remains safe but
retains more routes.
\end{remark}

\begin{remark}[Computation]
\label{rem:computation}
$|D(r)|\le 2^Bf$, where $f$ bounds the number of routes per bundle. Since
$\mu_r$ is concave, membership in $\Kstar_i$ is decided by evaluating $\mu_r$
at the endpoints of $\Theta$ and at the breakpoints of $E_r$, in
$O(|D(r)|\log|D(r)|)$ time; the whole frontier of a driver costs
$O(|R_i|\,2^Bf\log(2^Bf))$. The frontier is a post-processing step: it does
not speed up enumeration, but it shrinks every problem solved by
Algorithms~B and~C. Procedure~\ref{alg:frontier} states the computation.
\end{remark}

\begin{algorithm}[!htbp]
\caption{\textsc{Frontier}$(R_i,q,\Theta)$: the rule $\Pi^\star$ for one driver}
\label{alg:frontier}
\begin{algorithmic}[1]
\Require pool $R_i$ with bundles $S_r$ and coefficients $(K_r,W_r)$; FD prices $q$; $\Theta=[\lth,\uth]$
\Ensure local frontier $\Kstar_i\subseteq R_i$
\State $\Kstar_i\gets\emptyset$
\For{each $r\in R_i$}
  \State $D(r)\gets\{r'\in R_i\cup\{0_i\}\setminus\{r\}:S_{r'}\subseteq S_r\}$
  \State for each $r'\in D(r)$ form the line $g_{r'}(b)=K_{r'}+q(S_r\setminus S_{r'})+b\,W_{r'}$
  \State compute the lower envelope $E_r=\min_{r'\in D(r)}g_{r'}$ on $\Theta$ \Comment{$O(|D(r)|\log|D(r)|)$}
  \State $\mathcal B\gets\{\lth,\uth\}\cup\{\text{breakpoints of }E_r\text{ in }\Theta\}$
  \If{$\max_{b\in\mathcal B}\bigl(E_r(b)-K_r-b\,W_r\bigr)>0$} \Comment{$\mu_r$ is concave}
    \State add $r$ to $\Kstar_i$
  \EndIf
\EndFor
\State \Return $\Kstar_i$
\end{algorithmic}
\end{algorithm}

\begin{remark}[Relation to column dominance and redundancy]
\label{rem:lp}
In the linear relaxation of~\eqref{eq:wdp} let $\pi_o$ be the duals of the
order constraints and $\sigma_i\le 0$ that of driver $i$'s constraint. The FD
columns impose $\pi_o\le q_o$. If $S_{r'}\subseteq S_r$ and
$c_r\ge c_{r'}+q(S_r\setminus S_{r'})$, then, since both routes belong to
driver $i$, the term $\sigma_i$ cancels and
\[
\bar c_r-\bar c_{r'}=c_r-c_{r'}-\pi(S_r\setminus S_{r'})
\ge\sum_{o\in S_r\setminus S_{r'}}(q_o-\pi_o)\ge 0
\]
for every dual vector with $\pi\le q$. The rule $\Pi^\star$ thus deletes a
route when, at each admissible report, a route of the same driver dominates
it in reduced cost over the dual domain restricted by the outside option.
Equivalently, $r$ is redundant in the sense of \citet{sol1994column}, with
the combination consisting of $r'$ and the FD columns of
$S_r\setminus S_{r'}$.
\end{remark}

% ---------------------------------------------------------------------
\subsection{Payment computation}
\label{sec:sd-payments}

Exact payments~\eqref{eq:payment} require one removal solve per winner.
The reference implementation re-solves the base model once per winner with
all variables of that winner fixed to zero (Procedure~\ref{alg:payments}).

\begin{algorithm}[!htbp]
\caption{\textsc{Payments}$(R,b,x^\ast)$: exact Clarke-pivot payments (Algorithm C)}
\label{alg:payments}
\begin{algorithmic}[1]
\Require pools $R$ (fixed before reports), reports $b$, optimal allocation $x^\ast$ with value $V(R;b)$
\Ensure payments $(p_i)_{i\in N}$
\For{each driver $i$}
  \If{$i$ receives no route in $x^\ast$} $p_i\gets 0$; \textbf{continue} \EndIf
  \State fix $x_{ir}=0$ for all $r\in R_i$; re-solve~\eqref{eq:wdp} to optimality with the same tie-breaking
  \State $p_i\gets c_{ir_i^\ast}(b_i)+V_{-i}(R;b)-V(R;b)$
\EndFor
\State \Return $(p_i)_{i\in N}$
\end{algorithmic}
\end{algorithm}

% ---------------------------------------------------------------------
\subsection{Further results planned for the final thesis}
\label{sec:sd-planned}

This proposal establishes the safety of the frontier. Three further
questions from Section~\ref{sec:problem} will be developed, with statements
and proofs, in the final thesis. Preliminary checks for each are already
available.

\subsubsection{Minimality and worst-case size of the frontier (TQ2)}
\label{sec:sd-pol}

The plan is to show that the frontier cannot be improved by any rule of the
same kind: every route in $\Kstar_i$ should be indispensable in some instance
that looks identical from driver $i$'s own data, so that no safe local rule
can delete it. The planned proof adds to driver $i$ cheap single-order
competitors for every order outside the route's bundle. A second planned
result is a family of instances in which a gigworker's frontier contains
$\sum_{k\le B}\binom nk=\Theta(n^B)$ routes, none of which is ever optimal.
Such a result would show that the gap between gigworkers and occasional
drivers is structural. In a synthetic audit, every one of 1,157 frontier
routes was indeed indispensable in the constructed instance, and the planned
worst-case family behaved as predicted for $n=3,5,7,9$.

\subsubsection{Pruning during route generation (TQ3)}
\label{sec:sd-label}

The frontier is computed only after every route has been enumerated. A
\emph{label rule} would apply the same reasoning to partial routes during
enumeration: if dropping an order already picked up, and sending it to FD,
is guaranteed to be cheaper for every completion and every admissible report,
the partial route and all its extensions can be discarded. The direct
transfer of rollback pruning \citep{lozano2016exact} is unsafe here, because
working time is priced and waiting is allowed: a time saving on the known
part of the route can be absorbed by waiting later. A corrected rule with
\emph{junction} and \emph{absorption} terms has been designed and audited.
In an experimental build it left the frontier and every payment unchanged
(maximum difference $5.7\times10^{-14}$) and removed 15--43\% of label
extensions, but in pure Python its bookkeeping made route generation slower
(median time ratio 0.52--0.87). The final thesis will state the rule and
prove that it never changes the frontier.

\subsubsection{Decomposing the payment computation (TQ4)}
\label{sec:sd-decomp}

The \emph{conflict graph} links two drivers whenever some of their routes
share an order; it is built once from the pools, before any report is
read. When the FD price of each order is a fixed constant, the
counterfactual problem of a winner should involve only the connected
component of the graph that contains that winner, so each payment would
require re-solving a smaller problem. This property fails if the fixed fleet
has a shared capacity: removing one driver can then change the allocation in
a component that does not contain that driver. A small counterexample with
two orders and a fleet that can serve only one order has already been
constructed. On the experimental instances, the conflict graph was almost
always connected where the crowd is competitive, so the decomposition is
expected to bring little benefit there. The final thesis will give the
decomposition result with its proof.
% ---------------------------------------------------------------------
\subsection{A small numerical solution}
\label{sec:sd-example}

This section runs the complete procedure by hand on a small instance with
four orders, two gigworkers and two occasional drivers. All numbers are
illustrative; they were chosen so that every step can be checked by hand,
and every value below was recomputed by exhaustive enumeration.

\paragraph{Input.} The bundle cap is $B=2$, the distance coefficient is
$\kappa=0.5$ for every driver, the speed is 20\,km/h, every driver has
capacity 2, and the report domain is $\Theta=[18,25]$.
Table~\ref{tab:ex-input} lists the orders and drivers. The true values of
time are private; the platform never sees them.

\begin{table}[t]
\centering
\caption{Input of the numerical example}
\label{tab:ex-input}
\small
\setlength{\tabcolsep}{4pt}
\begin{tabular}{lccc@{\hspace{0.8cm}}llccc}
\toprule
\multicolumn{4}{c}{Orders} & \multicolumn{5}{c}{Drivers}\\
\cmidrule(r{0.6cm}){1-4}\cmidrule(l){5-9}
Order & Pickup & Delivery & $q_o$ & Driver & Class & Direct trip & Detour & True $\theta_i$\\
\midrule
$o_1$ & $P_1$ & $D_1$ & 15 & $g_1$ & GW & -- & -- & 20\\
$o_2$ & $P_2$ & $D_2$ & 45 & $g_2$ & GW & -- & -- & 24\\
$o_3$ & $P_3$ & $D_3$ & 40 & $c_1$ & OD & 10 km / 30 min & 20 min & 18\\
$o_4$ & $P_4$ & $D_4$ & 50 & $c_2$ & OD & 8 km / 24 min & 18 min & 22\\
\bottomrule
\end{tabular}
\end{table}

\paragraph{Step 1: route generation (Algorithm A).} For the occasional
driver $c_1$ the detour budget of 20 minutes eliminates every bundle that
contains $o_1$ or $o_4$, since a lower bound on their detour (34 and 29
minutes) already exceeds it; only $\{o_2\}$, $\{o_3\}$ and $\{o_2,o_3\}$
remain, out of ten bundles of size at most two. For the gigworker $g_1$ no
such bound exists: order $o_3$ is infeasible because of its time window, and
the pairs $\{o_1,o_4\}$ and $\{o_2,o_4\}$ are infeasible although each of
their orders is feasible on its own, which can only be discovered by
building the sequences. To keep the example small, it uses the nine routes
of Table~\ref{tab:ex-pool}; the routes of $g_1$ serving $\{o_4\}$ and of
$c_1$ serving $\{o_3\}$ alone are left out, and all values below refer to
this pool.

\begin{table}[t]
\centering
\caption{Route pool of the numerical example and reported costs at the true values of time}
\label{tab:ex-pool}
\small
\begin{tabular}{lllccccc}
\toprule
Route & Driver & Bundle & Distance & Time & $K$ & $W$ (h) & $c=K+\theta W$\\
\midrule
$r_1$ & $g_1$ & $\{o_1\}$ & 12 km & 36 min & 6.0 & 0.600 & 18.00\\
$r_2$ & $g_1$ & $\{o_2\}$ & 16 km & 48 min & 8.0 & 0.800 & 24.00\\
$r_3$ & $g_1$ & $\{o_1,o_2\}$ & 22 km & 72 min & 11.0 & 1.200 & 35.00\\
$r_4$ & $g_2$ & $\{o_3\}$ & 10 km & 30 min & 5.0 & 0.500 & 17.00\\
$r_5$ & $g_2$ & $\{o_4\}$ & 14 km & 42 min & 7.0 & 0.700 & 23.80\\
$r_6$ & $g_2$ & $\{o_3,o_4\}$ & 20 km & 60 min & 10.0 & 1.000 & 34.00\\
$r_7$ & $c_1$ & $\{o_2\}$ & 4 km (detour) & 15 min (detour) & 2.0 & 0.250 & 6.50\\
$r_8$ & $c_1$ & $\{o_2,o_3\}$ & 6 km (detour) & 20 min (detour) & 3.0 & 0.333 & 9.00\\
$r_9$ & $c_2$ & $\{o_4\}$ & 3 km (detour) & 15 min (detour) & 1.5 & 0.250 & 7.00\\
\bottomrule
\end{tabular}
\end{table}

\paragraph{Step 2: the local frontier.} For every route, the frontier
compares its cost with the envelope $E_r$ of~\eqref{eq:envelope} at the ends of
$\Theta$ (all envelopes here are linear on $\Theta$, so the endpoints
suffice). Table~\ref{tab:ex-frontier} gives the result. Only $r_1$ is
deleted: serving $o_1$ costs $g_1$ at least $16.80$ at any admissible report,
more than the FD price of 15. Route $r_3$ is kept although $r_1$ is deleted,
because the bundle $\{o_1,o_2\}$ is cheaper than $r_2$ plus FD for $o_1$
(margin 4.80 at $b=18$). This is the twin situation of Example~\ref{ex:twin}.

\begin{table}[t]
\centering
\caption{Local frontier of the numerical example, $\Theta=[18,25]$}
\label{tab:ex-frontier}
\small
\begin{tabular}{llcccccc}
\toprule
Route & $D(r)$ & $E_r(18)$ & $c_r(18)$ & $\mu_r(18)$ & $E_r(25)$ & $c_r(25)$ & In $\Kstar$?\\
\midrule
$r_1$ & $0$ & 15.00 & 16.80 & $-1.80$ & 15.00 & 21.00 & No\\
$r_2$ & $0$ & 45.00 & 22.40 & 22.60 & 45.00 & 28.00 & Yes\\
$r_3$ & $r_1,r_2,0$ & 37.40 & 32.60 & 4.80 & 43.00 & 41.00 & Yes\\
$r_4$ & $0$ & 40.00 & 14.00 & 26.00 & 40.00 & 17.50 & Yes\\
$r_5$ & $0$ & 50.00 & 19.60 & 30.40 & 50.00 & 24.50 & Yes\\
$r_6$ & $r_4,r_5,0$ & 59.60 & 28.00 & 31.60 & 64.50 & 35.00 & Yes\\
$r_7$ & $0$ & 45.00 & 6.50 & 38.50 & 45.00 & 8.25 & Yes\\
$r_8$ & $r_7,0$ & 46.50 & 9.00 & 37.50 & 48.25 & 11.33 & Yes\\
$r_9$ & $0$ & 50.00 & 6.00 & 44.00 & 50.00 & 7.75 & Yes\\
\bottomrule
\end{tabular}
\par\smallskip\footnotesize $0$ denotes the idle route; $E_r(b)$ for $r_3$ is attained by $r_2$ plus FD for $o_1$, and for $r_6$ by $r_5$ plus FD for $o_3$.
\end{table}

\paragraph{Step 3: winner determination (Algorithm B).} With truthful
reports, the optimal allocation gives $r_8=\{o_2,o_3\}$ to $c_1$ and
$r_9=\{o_4\}$ to $c_2$, and sends $o_1$ to FD:
\[
V=9.00+7.00+15=31.00 .
\]
The second-best allocation costs 34.00: it adds $r_1$, so that $g_1$ serves
$o_1$ at a reported cost of 18.00 instead of the FD price of 15. This is the
route deleted by the frontier, and it is used only by a worse allocation.
Both gigworkers lose. The same optimum is obtained on the pruned pool, as
Theorem~\ref{thm:frontier}(a) guarantees.

\paragraph{Step 4: payments (Algorithm C).} Each winner is removed in turn
and the problem re-solved:
\begin{align*}
V_{-c_1}&=c_{r_3}+c_{r_4}+c_{r_9}=35.00+17.00+7.00=59.00,\\
p_{c_1}&=c_{r_8}+V_{-c_1}-V=9.00+59.00-31.00=37.00,\\
V_{-c_2}&=q_{o_1}+c_{r_8}+c_{r_5}=15+9.00+23.80=47.80,\\
p_{c_2}&=c_{r_9}+V_{-c_2}-V=7.00+47.80-31.00=23.80.
\end{align*}
Without $c_1$, gigworker $g_1$ takes the bundle $\{o_1,o_2\}$ and $g_2$ takes
$o_3$; without $c_2$, gigworker $g_2$ takes $o_4$. Table~\ref{tab:ex-money}
summarises the financial outcome. The platform pays 75.80 in total, against
150 if every order were sent to FD; the difference between payments and true
costs, 44.80, is the information rent that buys truthfulness. All values,
including the counterfactuals, are identical on the pruned pool.

\begin{table}[t]
\centering
\caption{Financial outcome of the numerical example}
\label{tab:ex-money}
\small
\begin{tabular}{lccccc}
\toprule
Driver & Route & True cost & $V_{-i}$ & Payment $p_i$ & Utility $p_i-C_i$\\
\midrule
$c_1$ & $r_8$ & 9.00 & 59.00 & 37.00 & 28.00\\
$c_2$ & $r_9$ & 7.00 & 47.80 & 23.80 & 16.80\\
$g_1$, $g_2$ & -- & -- & -- & 0 & 0\\
\midrule
\multicolumn{2}{l}{Optimal value $V$} & \multicolumn{4}{l}{31.00 (routes 16.00, FD 15.00)}\\
\multicolumn{2}{l}{Total platform outlay} & \multicolumn{4}{l}{$37.00+23.80+15=75.80$ (all-FD: 150.00)}\\
\multicolumn{2}{l}{Information rent} & \multicolumn{4}{l}{$28.00+16.80=44.80$}\\
\bottomrule
\end{tabular}
\end{table}

\paragraph{Step 5: checking truthfulness.} Table~\ref{tab:ex-dsic} lets each
driver misreport within $\Theta$. A winner who misreports keeps the same
bundle, her payment does not change, and her utility does not change: the
Clarke-pivot payment depends on her own report only through which bundle she
wins. A loser who underbids to the lower end of $\Theta$ still loses. No
driver gains by misreporting, as Proposition~\ref{prop:dsic} guarantees.

\begin{table}[t]
\centering
\caption{Outcome of unilateral misreports in the numerical example}
\label{tab:ex-dsic}
\small
\begin{tabular}{lcccccc}
\toprule
Driver & True $\theta_i$ & Report $b_i$ & Wins & $V$ & Payment & Utility (true)\\
\midrule
$c_1$ & 18 & 18 (truthful) & $r_8$ & 31.00 & 37.00 & 28.00\\
$c_1$ & 18 & 18.5 & $r_8$ & 31.17 & 37.00 & 28.00\\
$c_1$ & 18 & 25 & $r_8$ & 33.33 & 37.00 & 28.00\\
$c_2$ & 22 & 22 (truthful) & $r_9$ & 31.00 & 23.80 & 16.80\\
$c_2$ & 22 & 18 & $r_9$ & 30.00 & 23.80 & 16.80\\
$c_2$ & 22 & 25 & $r_9$ & 31.75 & 23.80 & 16.80\\
$g_1$ & 20 & 18 & -- & 31.00 & 0 & 0\\
$g_2$ & 24 & 18 & -- & 31.00 & 0 & 0\\
\bottomrule
\end{tabular}
\end{table}

\paragraph{Step 6: payment decomposition.} Driver $g_1$ shares $o_2$ with
$c_1$, $c_1$ shares $o_3$ with $g_2$, and $g_2$ shares $o_4$ with $c_2$. The
conflict graph is a single component, so the planned decomposition of
Section~\ref{sec:sd-decomp} would give no saving here; this mirrors the finding of the computational study that the
graph is almost always connected where the crowd is competitive.

% ---------------------------------------------------------------------
\subsection{Implementation}
\label{sec:sd-impl}

The pipeline is implemented in Python~3.7.7. Table~\ref{tab:modules} maps the
components of Procedure~\ref{alg:mechanism} to the modules of the code base.

\begin{table}[t]
\centering
\caption{Implementation of the components}
\label{tab:modules}
\small
\begin{tabular}{P{3.6cm}P{6.0cm}P{4.6cm}}
\toprule
Component & Module (function) & Remarks\\
\midrule
Instance generator & \texttt{instance\_gen.py} (\texttt{generate\_instance}) & Locked, hashed parameters\\
Algorithm A & \texttt{dp\_labeling.py} (\texttt{build\_route\_pool}) $\to$ \texttt{t6\_dp.py} (\texttt{run\_dp}) & Layers 1 and 2\\
Frontier $\Pi^\star$ & \texttt{kstar\_rule.py} (\texttt{local\_frontier}) & Breakpoints of the lower envelope\\
Label rule (experimental) & \texttt{fdrule\_dp.py}; enumerator \texttt{dp\_fast.py} & Not used in the main runs\\
Algorithms B and C & \texttt{rq1\_wdp.py}, \texttt{rq\_common.py} (\texttt{vcg}) & CPLEX 12.10, one thread, zero gap\\
Conflict graph & \texttt{conflict\_graph.py} & Payment decomposition\\
Test oracle & \texttt{brute\_force.py} (\texttt{brute\_pool\_for\_driver}) & Exhaustive enumeration and WDP\\
\bottomrule
\end{tabular}
\end{table}

\paragraph{Solver settings and tie-breaking.} Algorithms~B and~C call CPLEX
12.10 through its Python interface with one thread and relative and absolute
gap tolerances of zero; any solve not reported as optimal is treated as a
failure and is not used. Tie-breaking is implemented in three steps: the
objective is fixed at its optimal value, the number of orders sent to FD is
minimised, and remaining ties are broken lexicographically by route
identifier. Algorithm~C uses the same pool and the same tie-breaking for
every counterfactual problem.

\paragraph{Bid independence.} Algorithm~A and the frontier take no report
as an argument. A hash of every driver's pool is computed before reports are
read, and a regression test checks that the pool and its hash are unchanged
when the reports change.

\paragraph{Verification.} Before any main run, every component was checked
against the independent oracle on instances with $n\in\{3,\dots,6\}$, two
time-window widths, three seeds and four drivers: 360 preliminary checks and
1,035 checks for the later research questions, all passed with deviations
below $10^{-6}$. The theoretical results were additionally audited with a
reference implementation that shares no code with the production one
(Table~\ref{tab:audit1}). It enumerates every precedence-feasible sequence,
computes the frontier analytically, and solves every winner-determination
problem exactly with the HiGHS solver on synthetic instances with $n=8$
orders, two gigworkers and two occasional drivers, $B=3$ and nine seeds.
Deliberate mutations (for example removing the FD term from the envelope)
were detected.

\begin{table}[t]
\centering
\caption{Synthetic audit of the local pruning frontier}
\label{tab:audit1}
\small
\begin{tabular}{P{10.2cm}P{4.0cm}}
\toprule
Check & Result\\
\midrule
Routes after Pareto filtering / retained by $\Pi^\star$ & 2,985 / 1,157 (38.8\%)\\
Max.\ $|V(R^{\Pi^\star})-V(R)|$ over base and removal solves & $5.7\times10^{-14}$\\
Max.\ payment difference (1,737 payments) & $5.7\times10^{-14}$\\
Routes optimal for some of 400 report profiles (three seeds) & 1.9--3.8\% of pool, all in $\Kstar$\\
Frontier routes deleted by a per-driver hull across bundles & 1,144 / 1,157\\
Mutation: FD term removed from $E_r$ & Detected\\
\bottomrule
\end{tabular}
\end{table}

% ---------------------------------------------------------------------
\subsection{Investigation of alternative solutions and the improved solution}
\label{sec:sd-alternatives}

The final design is the result of an iterative development in which each
candidate was tested against brute force and either rejected or improved.
Table~\ref{tab:history} summarises the development path.

\begin{table}[t]
\centering
\caption{Development path of the solution}
\label{tab:history}
\small
\begin{tabular}{P{0.8cm}P{4.0cm}P{5.0cm}P{3.6cm}}
\toprule
Step & Candidate & Finding & Decision\\
\midrule
1 & Level-wise insertion & Complete (0 mismatches, 1,152 instances); sequence count explodes under wide windows & Keep as baseline; needs compression\\
2 & Pareto dominance between sequences & 31.26\% of gigworker bundles lose sequences & Rejected\\
3 & Lookahead-profile dominance & Sound at depth 1; 66.5--72.4\% violations at depth $\ge 2$ & Rejected for generation\\
4 & Self-survival bound & Sound; prunes 0\% under wide windows & Rejected (no effect)\\
5 & Forward label-setting, exact key & 0 violations, including adversarial 2D cases; 5--10$\times$ faster & \textbf{Adopted as Algorithm A}\\
6 & Per-driver hull across bundles & Deletes routes that some instance needs (Example~\ref{ex:twin}) & Rejected\\
7 & Local FD-completion frontier & Safe (Theorem~\ref{thm:frontier}) & \textbf{Adopted as $\Pi^\star$}\\
8 & Rollback pruning transferred directly & Unsafe when waiting is priced & Corrected\\
9 & Label rule with junction and absorption terms & Frontier unchanged in audit; slower in pure Python & Planned (Section~\ref{sec:sd-label})\\
10 & Component decomposition of payments & Exact under separability; fails with shared FD capacity & Planned (Section~\ref{sec:sd-decomp})\\
\bottomrule
\end{tabular}
\end{table}

The improvement from step~1 to step~5 is in completeness and speed: the
label-setting algorithm is provably complete and 5--10 times faster than the
level-wise baseline. The improvement from step~6 to step~7 is in safety:
the per-driver hull is strong but wrong, whereas the frontier is exact for
every admissible report (Theorem~\ref{thm:frontier}). Whether any other
local rule can do better is the subject of the minimality result planned in
Section~\ref{sec:sd-pol}.

% ---------------------------------------------------------------------
\subsection{Comparison of alternative solutions under feasibility constraints and standards}
\label{sec:sd-compare}

\paragraph{Economic feasibility.} On the numerical example,
Table~\ref{tab:ex-alt} compares the selected design with alternative
designs. Using both crowd classes is as good as the better single class here
(both 31.00), and much better than gigworkers alone (69.00). Allowing bundles
lowers the cost from 45.50 (single-order routes) to 31.00, a 31.9\%
reduction. A posted price of $0.8\,q_o$ attracts a cover of all orders, but
since the platform cannot see costs, its payout (120.00) is higher than the
VCG outlay (75.80) and the selected routes have a true cost of at least
59.00. This small example is illustrative only: in the computational study
the optimised posted price paid less than VCG on average but lost about 4\%
efficiency, and the value of joint bidding and bundling depended on whether
the crowd could compete with the fixed fleet (Chapter~5).

\begin{table}[t]
\centering
\caption{Alternative designs on the numerical example (true values of time)}
\label{tab:ex-alt}
\small
\setlength{\tabcolsep}{4pt}
\begin{tabular}{P{4.4cm}P{1.9cm}P{1.9cm}P{2.0cm}P{3.6cm}}
\toprule
Design & True system cost & Platform outlay & Truthful? & Selected routes\\
\midrule
All orders to FD & 150.00 & 150.00 & -- & none\\
Gigworkers only (VCG) & 69.00 & -- & Yes & $r_3$, $r_6$\\
Occasional drivers only (VCG) & 31.00 & -- & Yes & $r_8$, $r_9$\\
Single-order routes only (VCG) & 45.50 & -- & Yes & $r_4$, $r_7$, $r_9$\\
Posted price $0.8\,q_o$ & $\ge 59.00$ & 120.00 & Accept or decline only & e.g.\ $r_3$, $r_4$, $r_9$\\
\textbf{Joint bidding, bundles, VCG (selected)} & \textbf{31.00} & \textbf{75.80} & \textbf{Yes} & $r_8$, $r_9$, FD for $o_1$\\
\bottomrule
\end{tabular}
\par\smallskip\footnotesize Outlay is reported only where a payment rule is specified for the design. Under the posted price every cover of all orders has the same payout, so the true cost depends on tie-breaking.
\end{table}

\paragraph{Environmental aspects.} The mechanism influences emissions
through the distance driven. Bundling lets one driver serve several orders on
one route, and occasional drivers are charged only for their detour, so
orders that lie along their trips are served with little extra driving. In
the computational study, larger bundles reduced the share of orders sent to
the fixed fleet from 0.897 (single-order routes) to 0.731 with $B=3$. The
model does not price emissions; the ISO~14083:2023 standard for quantifying
greenhouse-gas emissions of transport chains would provide a consistent basis
for converting the distances of selected routes into emissions in an
extension.

\paragraph{Social aspects.} For drivers, the mechanism makes truthful
reporting a dominant strategy and guarantees nonnegative utility to
winners, so drivers with little experience in bidding are not disadvantaged
by others' strategic skill. Publishing a hash of the route pool before
bidding lets drivers verify that the allocation range was not changed after
the bids were seen. The pre-auction certificate avoids inviting drivers who
cannot win at any admissible report. The mechanism does not address
minimum earnings, working-hour limits or fairness across drivers, which
platforms must handle under applicable labour regulation.

\paragraph{Data and information security.} The mechanism collects one
number per driver and uses location and trip data of occasional drivers.
Such data are personal data; a deployment in Vietnam would have to comply
with the Law on Personal Data Protection (Law No.~91/2025/QH15, in force
since 1 January 2026) and its guiding Decree No.~356/2025/ND-CP, which
replaced Decree No.~13/2023/ND-CP, and an information-security management
system such as ISO/IEC~27001 would govern their storage and access.
Standards such as OSHA, HACCP or TCVN product standards do not apply directly
to a procurement algorithm.

\paragraph{Summary.} Among the alternatives, the selected design is the only
one that satisfies all correctness requirements R1--R7 simultaneously:
exact VCG on a bid-independent pool reduced to the local frontier. Its costs
are a payout premium over true costs (the information rent) and the work of
enumerating the pool, which for gigworkers grows quickly with the number of
orders. Chapter~5 quantifies both on
the full experimental grid.
